\section{格点正规化下非局域算符的重正化}

式~(\ref{Eq:lag}) 中的 HQET 拉氏量取了重夸克质量无穷大的极限，因而不包含任何质量项。重夸克的质量修正只能接受幂次发散的贡献~\cite{Maiani:1991az}，而这种贡献在维数正规化中并不存在。

然而，在诸如格点正规化这类紫外截断正规化中，当超出领头阶微扰论时，重夸克的自能会引入线性发散，它必须被吸收进一个有效的质量抵消项，
\begin{equation}
    \delta {\cal L}_m =  -\delta m\overline{Q} Q
\end{equation}
其中 $\delta m$ 为 ${\cal O}(1/a)$ 量级，$a$ 为格距~\cite{Maiani:1991az}。尽管格点正规化破坏洛伦兹对称性并引入算符混合，但这对 $O(z_2,z_1)$ 并不构成问题，因为在格点 QCD 中不存在能与它混合的更低维算符。
在 HQET 框架中，线性发散的出现很容易理解：既然重夸克无限重，它的行为就像一个静态色源，其能量将具有类库仑的形式 $1/r$，因此当色源是点粒子时能量会线性发散。在我们的辅助重夸克拉氏量中，作为静态色源的物理图像已不复存在，但线性发散的去除基本上可以用同样的方式完成。由于线性发散只出现在辅助重夸克的自能之中，它可以像文献~\cite{Maiani:1991az} 那样被吸收进重夸克质量的重正化。

于是，总的重夸克拉氏量变为
\begin{equation}\label{Eq:spaceliken}
    {\cal L}_Q = \bar Q (i n\cdot D-\delta m)Q.
\end{equation}
注意，拉氏量的 BRST 不变性要求 $\delta m$ 依赖于 $n$ 的符号~\cite{Dorn:1986dt}，这使得对类光的 $n^\mu$ 有 $\delta m=0$。在类空 $n^\mu$ 的情形，式~(\ref{Eq:spaceliken}) 中的 $\delta m$ 是虚数，因此我们可以写 $\delta m=i\delta {\bar m}$。在对重场积分之后，重正化的算符变为
\begin{equation}\label{Eq:cutoffrenorm}
        O_R = Z_{\bar j}^{-1}Z_{j}^{-1}e^{\delta \bar m|z_2-z_1|}\overline{\psi}(z_2) \Gamma L(z_2,z_1)\psi(z_1).
\end{equation}
这样，指数中出现了一个依赖 $z$ 的质量抵消项。
质量重正化之后，至多剩下对数发散，它们将被来自拉氏量的抵消项或轻重复合算符重正化中的抵消项所消去。上述重正化形式此前已在多篇文献中被提出~\cite{Musch:2010ka,Ishikawa:2016znu,Chen:2016fxx,Constantinou:2017sej,Alexandrou:2017huk,Chen:2017mzz}，
但未被证明。

值得指出的是，格点正规化下的显式单圈计算~\cite{Chen:2016fxx} 确实表明，式~(\ref{Eq:cutoffrenorm}) 形式的指数质量重正化，在傅里叶变换到动量空间后恰好消去了文献~\cite{Xiong:2013bka} 中准 PDF 的线性发散，其中
\beq
\delta \bar m=\frac{2\pi\alpha_s }{3a}.
\eeq

当 $z_2<z_1$ 时，上述重正化同样适用。
